\documentclass[11pt,a4paper]{article}
\usepackage{amsmath,amssymb,amsthm}
\usepackage{geometry}
\geometry{margin=1in}
\usepackage{hyperref}

\newtheorem{theorem}{Theorem}[section]
\newtheorem{lemma}[theorem]{Lemma}
\newtheorem{proposition}[theorem]{Proposition}
\newtheorem{definition}[theorem]{Definition}

\title{\bfseries Harmonic Sine Transform Programme II\\
\large Greedy Harmonic Tree and Orthonormal Haar Basis}
\author{}
\date{}

\begin{document}
\maketitle

\begin{abstract}
Building on the greedy harmonic expansion of angles established in Programme~I, we construct the associated binary tree of cylinder intervals—the greedy harmonic tree—and extract an orthonormal Haar basis \(\mathcal{B}\) of the Hilbert space \(L^2[0,\pi]\).  We prove that \(\mathcal{B}\) is a complete orthonormal system, thereby defining the Harmonic Sine Transform (HST) as an isometric isomorphism onto \(\ell^2(\mathcal{B})\).  This basis is the central tool for all subsequent analytic and number‑theoretic developments.
\end{abstract}

\section{Introduction}
The greedy harmonic expansion
\[
\frac{\theta}{\pi} = \sum_{n=0}^{\infty} \frac{\delta_n(\theta)}{n+2},\qquad \theta\in[0,\pi],
\]
with binary digits \(\delta_n(\theta)\in\{0,1\}\) determined by the threshold rule, induces a partition of the interval \([0,\pi]\) into cylinder sets.  These sets form the nodes of a binary tree \(\mathcal{T}\) whose levels correspond to successive digits.  The present programme constructs an orthonormal basis of Haar type from this tree.

\section{Goals of the Programme}
The programme aims to prove the following statements.
\begin{enumerate}
\item The greedy harmonic tree \(\mathcal{T}\) is a well‑defined binary tree whose nodes at depth \(n\) are the cylinder intervals \(I_\varepsilon\) (\(\varepsilon\in\{0,1\}^n\)).
\item For each \emph{split} node (an interval whose length exceeds \(\pi/(n+2)\)), we define a Haar wavelet \(\psi_\varepsilon\in L^2[0,\pi]\) with unit norm and zero mean.
\item The set \(\mathcal{B}\) consisting of the constant function \(\psi_\varnothing = 1/\sqrt{\pi}\) and all split‑node wavelets is orthonormal.
\item \(\mathcal{B}\) spans a dense subspace of \(L^2[0,\pi]\); hence it is an orthonormal basis.
\item The Harmonic Sine Transform (HST) \(\mathcal{H}:L^2[0,\pi]\to\ell^2(\mathcal{B})\), given by \(\mathcal{H}f(\psi)=\langle f,\psi\rangle\), is an isometric isomorphism.
\end{enumerate}

\section{Background}
We recall from Programme~I the essential objects.  Let \(\alpha_n = \pi/(n+2)\) for \(n\ge 0\).  The threshold angles \(\theta_n^*\in[0,\pi]\) satisfy
\[
\delta_n(\theta) =
\begin{cases}
0, & \theta < \theta_n^*,\\
1, & \theta \ge \theta_n^*.
\end{cases}
\]
For a binary prefix \(\varepsilon = (\varepsilon_0,\dots,\varepsilon_{n-1})\in\{0,1\}^n\), the cylinder interval \(I_\varepsilon\) consists of all \(\theta\in[0,\pi]\) whose first \(n\) digits equal \(\varepsilon\).  We denote its length by \(\ell(\varepsilon)\).  The recursive length rule is
\[
\ell(\varepsilon1) = \alpha_n,\qquad \ell(\varepsilon0) = \ell(\varepsilon) - \alpha_n \quad\text{if }\ell(\varepsilon) > \alpha_n,
\]
and if \(\ell(\varepsilon) \le \alpha_n\) the node \(\varepsilon\) is \emph{unsplit}: only the left child exists and has the same length.

The collection of all \(I_\varepsilon\) (with the empty prefix giving the whole interval \([0,\pi]\)) forms the \textbf{greedy harmonic tree} \(\mathcal{T}\).  Every \(\theta\in[0,\pi]\) corresponds to a unique infinite path; the maximal elements (limit of intervals along a path) are singletons.

\section{Construction of the Haar Wavelets}
For a node \(\varepsilon\) of length \(n\) that \emph{splits} (i.e., \(\ell(\varepsilon) > \alpha_n\)), we have two children:
\[
I_{\varepsilon0} = [a, \theta_n^*),\qquad I_{\varepsilon1} = [\theta_n^*, b),
\]
with lengths \(\ell(\varepsilon0) = \ell(\varepsilon)-\alpha_n\) and \(\ell(\varepsilon1) = \alpha_n\).  Define the \textbf{Haar wavelet}
\begin{equation}\label{eq:wavelet}
\psi_\varepsilon(\theta) =
\frac{1}{\sqrt{\alpha_n}}\mathbf{1}_{I_{\varepsilon1}}(\theta) \;-\;
\frac{1}{\sqrt{\ell(\varepsilon)-\alpha_n}}\mathbf{1}_{I_{\varepsilon0}}(\theta).
\end{equation}
For the root (empty prefix) we use the constant function
\[
\psi_\varnothing(\theta) = \frac{1}{\sqrt{\pi}}.
\]
The set \(\mathcal{B}\) comprises \(\psi_\varnothing\) together with all \(\psi_\varepsilon\) for split nodes \(\varepsilon\).

\section{Orthonormality}
\begin{theorem}\label{thm:orthonormal}
The system \(\mathcal{B}\) is orthonormal in \(L^2[0,\pi]\).
\end{theorem}
\begin{proof}
The constant function has norm \(1\) because \(\int_0^\pi (1/\sqrt{\pi})^2\,d\theta = 1\).  For a split node \(\varepsilon\) of depth \(n\), the two children are disjoint and exhaust \(I_\varepsilon\).  Hence
\[
\int_0^\pi |\psi_\varepsilon|^2 = \frac{1}{\alpha_n}\int_{I_{\varepsilon1}}\!1 \;+\; \frac{1}{\ell(\varepsilon)-\alpha_n}\int_{I_{\varepsilon0}}\!1
= \frac{\alpha_n}{\alpha_n} + \frac{\ell(\varepsilon)-\alpha_n}{\ell(\varepsilon)-\alpha_n} = 2.
\]
Wait, careful: Actually the sum is \(\frac{1}{\alpha_n}\cdot\alpha_n + \frac{1}{\ell(\varepsilon)-\alpha_n}\cdot(\ell(\varepsilon)-\alpha_n) = 1+1 = 2\). That would give norm \(\sqrt{2}\), not 1. There is a mistake: We need to adjust coefficients. In a standard Haar basis, the wavelet on an interval is \(\frac{1}{\sqrt{|L|}}\mathbf{1}_L - \frac{1}{\sqrt{|R|}}\mathbf{1}_R\) but its norm squared is \(\frac{1}{|L|}|L| + \frac{1}{|R|}|R| = 1+1=2\). So we need to divide by \(\sqrt{2}\). But the standard Haar wavelet uses normalization such that the norm is 1, so we should use \(\frac{1}{\sqrt{2}}(\frac{1}{\sqrt{|L|}}\mathbf{1}_L - \frac{1}{\sqrt{|R|}}\mathbf{1}_R)\)? Actually the usual Haar wavelet on [0,1] with split at 1/2 is \(\psi = \mathbf{1}_{[0,1/2)} - \mathbf{1}_{[1/2,1)}\); its norm squared is 1/2+1/2=1, so it's normalized. For unequal splits, to have norm 1 we need to set \(\psi_\varepsilon = \frac{1}{\sqrt{\alpha_n (\ell(\varepsilon)-\alpha_n)/\ell(\varepsilon)}} (\dots)\). Wait, let's compute correctly.

We want \(\int \psi_\varepsilon^2 = 1\). If we set \(\psi_\varepsilon = a \mathbf{1}_{I_{\varepsilon1}} - b \mathbf{1}_{I_{\varepsilon0}}\), then the norm squared is \(a^2 |I_{\varepsilon1}| + b^2 |I_{\varepsilon0}| = 1\). Also we might want the wavelet to have zero mean: \(a|I_{\varepsilon1}| - b|I_{\varepsilon0}| = 0\). The zero mean condition forces \(b = a |I_{\varepsilon1}| / |I_{\varepsilon0}| = a \alpha_n/(\ell(\varepsilon)-\alpha_n)\). Then norm squared becomes \(a^2 \alpha_n + a^2 \frac{\alpha_n^2}{(\ell(\varepsilon)-\alpha_n)^2} (\ell(\varepsilon)-\alpha_n) = a^2 \alpha_n (1 + \frac{\alpha_n}{\ell(\varepsilon)-\alpha_n}) = a^2 \alpha_n \frac{\ell(\varepsilon)}{\ell(\varepsilon)-\alpha_n}\). Set this equal to 1 gives \(a = \sqrt{ \frac{\ell(\varepsilon)-\alpha_n}{\alpha_n \ell(\varepsilon)} }\). Then \(b = a \frac{\alpha_n}{\ell(\varepsilon)-\alpha_n} = \sqrt{ \frac{\alpha_n}{(\ell(\varepsilon)-\alpha_n) \ell(\varepsilon)} }\). So a possible normalized zero‑mean wavelet is
\[
\psi_\varepsilon = \frac{1}{\sqrt{\ell(\varepsilon)}} \left( \sqrt{ \frac{\ell(\varepsilon)-\alpha_n}{\alpha_n} } \mathbf{1}_{I_{\varepsilon1}} - \sqrt{ \frac{\alpha_n}{\ell(\varepsilon)-\alpha_n} } \mathbf{1}_{I_{\varepsilon0}} \right).
\]
This is a bit messy. Simpler: we can adopt wavelets that are not normalized to 1 directly, but then we need to renormalize the orthonormal basis. The Haar basis constructed from a general partition often uses normalized wavelets as in \cite{wavelets}. The version in the original HST paper used the simpler difference but then normalized by weights connected to lengths. Actually, the paper defined \(\psi_\varepsilon = \frac{1}{\sqrt{\alpha_n}}\mathbf{1}_{I_{\varepsilon1}} - \frac{1}{\sqrt{\ell(\varepsilon)-\alpha_n}}\mathbf{1}_{I_{\varepsilon0}}\), and claimed orthonormality? I recall that in the original "harmonic-tree-haar.tex" they might have defined differently. Let's check the document "harmonic-tree-haar.tex" mentioned in the research guide: they might have used a normalized version. In the abstract of the "Harmonic Sine Transform" paper, they said "For every split node ... we define a Haar wavelet (with unit L2 norm)". So they must have a proper normalization. I'll adopt the formula I derived.

Now orthogonality: wavelets at different nodes have disjoint support, because they are supported on the children of a node, and the children of distinct nodes at different depths are either disjoint or one is contained in the other. The wavelets are orthogonal to the constant function by zero mean. All good if we use normalized zero-mean wavelets. The set is orthonormal.

Thus we need to formalize the correct wavelet definition. I'll present it cleanly.

\end{proof}

We redefine the Haar wavelet properly:

For a split node \(\varepsilon\) with children \(L = I_{\varepsilon0}\), \(R = I_{\varepsilon1}\) and lengths \(|L| = \ell(\varepsilon) - \alpha_n\), \(|R| = \alpha_n\), define
\[
\psi_\varepsilon = \frac{1}{\sqrt{|L||R|/|I_\varepsilon|}} \left( \frac{1}{\sqrt{|R|}}\mathbf{1}_{R} - \frac{1}{\sqrt{|L|}}\mathbf{1}_{L} \right)
= \sqrt{ \frac{\ell(\varepsilon)}{\alpha_n (\ell(\varepsilon)-\alpha_n)} } \left( \sqrt{ \frac{\ell(\varepsilon)-\alpha_n}{\alpha_n} } \mathbf{1}_{R} - \sqrt{ \frac{\alpha_n}{\ell(\varepsilon)-\alpha_n} } \mathbf{1}_{L} \right).
\]
Simplify: 
\[
\psi_\varepsilon = \frac{1}{\sqrt{ \alpha_n (\ell(\varepsilon)-\alpha_n) / \ell(\varepsilon) }} \left( \frac{1}{\sqrt{\alpha_n}} \mathbf{1}_{R} - \frac{1}{\sqrt{\ell(\varepsilon)-\alpha_n}} \mathbf{1}_{L} \right) ? 
\]
Let's just define a clean version: set
\[
\psi_\varepsilon = \frac{1}{\sqrt{ \frac{\alpha_n(\ell(\varepsilon)-\alpha_n)}{\ell(\varepsilon)} }} \left( \frac{1}{\sqrt{\alpha_n}} \mathbf{1}_{R} - \frac{1}{\sqrt{\ell(\varepsilon)-\alpha_n}} \mathbf{1}_{L} \right).
\]
Then the factor out front is \(\sqrt{ \frac{\ell(\varepsilon)}{\alpha_n(\ell(\varepsilon)-\alpha_n)} }\). So the total coefficient on \(\mathbf{1}_R\) is \(\sqrt{ \frac{\ell(\varepsilon)}{\alpha_n(\ell(\varepsilon)-\alpha_n)} } \cdot \frac{1}{\sqrt{\alpha_n}} = \frac{\sqrt{\ell(\varepsilon)}}{\alpha_n\sqrt{\ell(\varepsilon)-\alpha_n}}\), and on \(\mathbf{1}_L\) is \(-\frac{\sqrt{\ell(\varepsilon)}}{\sqrt{\alpha_n}(\ell(\varepsilon)-\alpha_n)}\). Norm squared = \(\frac{\ell(\varepsilon)}{\alpha_n^2 (\ell(\varepsilon)-\alpha_n)} \cdot \alpha_n + \frac{\ell(\varepsilon)}{\alpha_n (\ell(\varepsilon)-\alpha_n)^2} \cdot (\ell(\varepsilon)-\alpha_n) = \frac{\ell(\varepsilon)}{\alpha_n (\ell(\varepsilon)-\alpha_n)} + \frac{\ell(\varepsilon)}{\alpha_n (\ell(\varepsilon)-\alpha_n)} = 2 \frac{\ell(\varepsilon)}{\alpha_n (\ell(\varepsilon)-\alpha_n)}\). Wait, that's not 1. I'm getting confused.

Let's compute properly from the general construction: We want a function supported on \(L\cup R\), piecewise constant, orthogonal to constants, unit norm. The space of such functions on \(L\cup R\) (relative to the parent interval) is one-dimensional. A convenient orthonormal basis vector is 
\[
\psi = \frac{1}{\sqrt{ \frac{|L||R|}{|L|+|R|} }} \left( \frac{\mathbf{1}_R}{|R|} - \frac{\mathbf{1}_L}{|L|} \right).
\]
Then \(\int \psi^2 = \frac{1}{ \frac{|L||R|}{|L|+|R|} } \left( \frac{|R|}{|R|^2} + \frac{|L|}{|L|^2} \right) = \frac{|L|+|R|}{|L||R|} \cdot \left( \frac{1}{|R|} + \frac{1}{|L|} \right) = \frac{|L|+|R|}{|L||R|} \cdot \frac{|L|+|R|}{|L||R|} = \left( \frac{|L|+|R|}{|L||R|} \right) \cdot (|L|+|R|) = \frac{(|L|+|R|)^2}{|L||R|}\). Not 1. So that's wrong.

The standard formula for a Haar wavelet on an interval split into two children is:
\[
\psi_{j,k} = \frac{1}{\sqrt{|L|}} \mathbf{1}_L - \frac{1}{\sqrt{|R|}} \mathbf{1}_R
\]
has norm \(\sqrt{2}\). To get norm 1, use \(\frac{1}{\sqrt{2}} \left( \frac{1}{\sqrt{|L|}} \mathbf{1}_L - \frac{1}{\sqrt{|R|}} \mathbf{1}_R \right)\). Then norm is 1. Zero mean: indeed, \(\int \psi = \frac{1}{\sqrt{2}} ( \frac{|L|}{\sqrt{|L|}} - \frac{|R|}{\sqrt{|R|}} ) = \frac{1}{\sqrt{2}} (\sqrt{|L|} - \sqrt{|R|}) \neq 0\) unless \(|L|=|R|\). So it's not zero mean.

Thus to have both unit norm and zero mean, we need the function to be proportional to \(\mathbf{1}_L/|L| - \mathbf{1}_R/|R|\) which has zero mean automatically. Then choose the constant to achieve norm 1. Compute \(f = c ( \frac{\mathbf{1}_L}{|L|} - \frac{\mathbf{1}_R}{|R|} )\). Then \(\|f\|^2 = c^2 ( \frac{1}{|L|^2}|L| + \frac{1}{|R|^2}|R| ) = c^2 ( \frac{1}{|L|} + \frac{1}{|R|} ) = c^2 \frac{|L|+|R|}{|L||R|}\). Set this equal to 1 gives \(c = \sqrt{ \frac{|L||R|}{|L|+|R|} }\). So the wavelet is
\[
\psi_\varepsilon = \sqrt{ \frac{|L||R|}{|L|+|R|} } \left( \frac{\mathbf{1}_L}{|L|} - \frac{\mathbf{1}_R}{|R|} \right)
= \sqrt{ \frac{\alpha_n(\ell(\varepsilon)-\alpha_n)}{\ell(\varepsilon)} } \left( \frac{\mathbf{1}_{I_{\varepsilon0}}}{\ell(\varepsilon)-\alpha_n} - \frac{\mathbf{1}_{I_{\varepsilon1}}}{\alpha_n} \right).
\]
This is zero mean and unit norm. Moreover, it's orthogonal to constants. This is the correct, rigorous wavelet.

Thus, in the construction we should define \(\psi_\varepsilon\) exactly as above.

Now, orthogonality between different wavelets: If two wavelets correspond to nodes where one node's interval is strictly contained in another but not equal, they are orthogonal because the constant value on each piece? Actually, the wavelets are supported on the parent's children; if two wavelets are from nodes where one interval is contained in a left child of another, then the inner product may not be zero automatically. However, standard Haar wavelet construction on a binary tree ensures orthonormality if we build them in this way (they become orthogonal due to the vanishing moments and the nesting). In fact, the construction above yields the orthonormal Haar basis for the binary partition. I recall the general construction of Haar basis on a dyadic tree: each wavelet is orthogonal to all others. The proof uses that each wavelet is constant on each of its two children, and its integral over the parent is zero. For two different nodes, either their supports are disjoint, or one is contained in a child of the other but then the inner product involves integration over a region where one wavelet is constant and the other integrates to zero. Let's check: Suppose \(\varepsilon\) is a split node and \(\eta\) is a node that contains \(\varepsilon\) properly. Then the wavelet \(\psi_\eta\) is constant on the children of \(\eta\), one of which contains \(\varepsilon\). But \(\psi_\varepsilon\) is supported inside one of those children and has zero mean over its parent, but not necessarily over that child. The inner product will be the constant value of \(\psi_\eta\) on that child times the integral of \(\psi_\varepsilon\) over that child. The integral of \(\psi_\varepsilon\) over its support is zero, but over a child of \(\eta\) that strictly contains \(\varepsilon\)'s parent, the integral of \(\psi_\varepsilon\) is zero because it's supported inside a subinterval and its integral over its own parent is zero, but the child of \(\eta\) may be larger than the parent of \(\varepsilon\), and the integral of \(\psi_\varepsilon\) over that larger interval might not be zero. So orthogonality fails with the simple zero-mean definition. The classical Haar basis on a dyadic tree works because the wavelet on a node is defined as the difference of normalized indicators of its two children, and then orthogonalized to all previous ones via Gram-Schmidt across the tree. That yields orthonormal basis. In that case, the wavelet for a node is not simply the zero-mean normalized one; it's the result of a specific construction that ensures orthogonality with all finer scales. 

Thus, the simple "Haar-like" wavelets defined locally may not be orthogonal across different levels unless the tree is dyadic and splits are symmetric. The greedy harmonic tree is highly asymmetric. To get an orthonormal basis, we need to use the **orthonormal Haar basis for an arbitrary binary tree partition** as constructed in wavelet theory using the so-called "Haar basis on a metric tree" or "binary tree orthonormal basis". This is a known construction: given a binary tree of intervals that partitions the space down to singletons, there is an orthonormal basis consisting of functions that are constant on the leaves of each subtree. It can be obtained by taking for each node a function that is the difference of the normalized characteristic functions of its two children, and then Gram-Schmidt orthogonalization over the tree. The result is exactly the basis described in the HST paper. So we need to present the rigorous construction.

I'll outline: start with the constant function. Then for each split node, define a function \(f_\varepsilon = \mathbf{1}_{I_{\varepsilon0}}/\ell(\varepsilon0) - \mathbf{1}_{I_{\varepsilon1}}/\ell(\varepsilon1)\). This is zero mean. Then perform Gram-Schmidt on the set of all \(f_\varepsilon\) (ordered by depth) to obtain an orthonormal set \(\psi_\varepsilon\). Then since the original functions span the same space as the indicator functions of cylinder intervals, which are dense, the orthonormalized set is complete. This is rigorous and standard.

Thus the programme will describe that construction and prove it yields an orthonormal basis. This may be a bit heavy but can be summarized.

I'll craft the document: define the initial set of Haar candidates \(f_\varepsilon\), then orthogonalize them to the constant and all previously obtained wavelets at smaller depths (i.e., larger intervals). Show that the resulting \(\psi_\varepsilon\) form an orthonormal basis. The orthogonality across different nodes is guaranteed by construction. Completeness follows because the linear span of the \(f_\varepsilon\) and constant is dense.

We also note that the constant function is \(\psi_\varnothing\). So the HST is defined as the map taking inner products with these \(\psi\)'s.

That will be Programme II.

Now output LaTeX.```latex
\documentclass[11pt,a4paper]{article}
\usepackage{amsmath,amssymb,amsthm}
\usepackage{geometry}
\geometry{margin=1in}
\usepackage{hyperref}

\newtheorem{theorem}{Theorem}[section]
\newtheorem{lemma}[theorem]{Lemma}
\newtheorem{proposition}[theorem]{Proposition}
\newtheorem{definition}[theorem]{Definition}

\title{\bfseries Harmonic Sine Transform Programme II\\
\large Greedy Harmonic Tree and Orthonormal Haar Basis}
\author{}
\date{}

\begin{document}
\maketitle

\begin{abstract}
Continuing from Programme~I, we construct the greedy harmonic tree \(\mathcal{T}\) and extract from it an orthonormal Haar basis \(\mathcal{B}\) of \(L^2[0,\pi]\).  The basis consists of the constant function and a family of piecewise‑constant wavelets attached to the split nodes of \(\mathcal{T}\).  Orthonormality is achieved by a natural Gram–Schmidt orthogonalisation of elementary zero‑mean functions.  Completeness follows because the original generating functions span all cylindrical step functions, which are dense in \(L^2[0,\pi]\).  The resulting Harmonic Sine Transform (HST) is an isometric isomorphism onto \(\ell^2(\mathcal{B})\).
\end{abstract}

\section{Introduction}
Programme~I established the greedy harmonic expansion
\[
\frac{\theta}{\pi} = \sum_{n=0}^{\infty} \frac{\delta_n(\theta)}{n+2},\qquad \theta\in[0,\pi],
\]
and the associated cylinder intervals \(I_\varepsilon\) indexed by binary prefixes \(\varepsilon\in\{0,1\}^n\).  These intervals form the nodes of a binary tree \(\mathcal{T}\)—the \emph{greedy harmonic tree}—whose levels are determined by the digits.  In the present programme we use this tree to construct an orthonormal basis of \(L^2[0,\pi]\) consisting of piecewise‑constant functions.  This basis will serve as the fundamental tool for the Harmonic Sine Transform (Programme~III) and all subsequent operator‑theoretic constructions.

\section{Goals of the Programme}
The programme aims to prove the following.
\begin{enumerate}
\item The greedy harmonic tree \(\mathcal{T}\) is a well‑defined binary tree whose nodes at depth \(n\) are the cylinder intervals \(I_\varepsilon\) of length \(\ell(\varepsilon)\).
\item For every \emph{split} node \(\varepsilon\) (where \(\ell(\varepsilon) > \pi/(n+2)\)) one can define an elementary zero‑mean function \(f_\varepsilon\).
\item Orthonormalising the family \(\{f_\varepsilon\}\) together with the constant function yields an orthonormal system \(\mathcal{B} = \{\psi_\varnothing\} \cup \{\psi_\varepsilon : \varepsilon \text{ split}\}\).
\item The system \(\mathcal{B}\) is complete in \(L^2[0,\pi]\); hence it is an orthonormal basis.
\item The Harmonic Sine Transform \(\mathcal{H}: L^2[0,\pi] \to \ell^2(\mathcal{B})\) is well‑defined and an isometric isomorphism.
\end{enumerate}

\section{Preliminaries}
We recall the notation from Programme~I.  Let \(\alpha_n = \pi/(n+2)\) for \(n\ge 0\).  The threshold angles \(\theta_n^*\) satisfy \(\delta_n(\theta)=1 \iff \theta \ge \theta_n^*\).  For a binary prefix \(\varepsilon = (\varepsilon_0,\dots,\varepsilon_{n-1})\in\{0,1\}^n\), the cylinder set
\[
I_\varepsilon = \bigl\{ \theta\in[0,\pi] : \delta_k(\theta)=\varepsilon_k,\; 0\le k<n \bigr\}
\]
is a half‑open interval (or singleton) of length \(\ell(\varepsilon)\).  The splitting rule is:
\[
\ell(\varepsilon1) = \alpha_n,\qquad
\ell(\varepsilon0) = \ell(\varepsilon) - \alpha_n \quad\text{if }\ell(\varepsilon) > \alpha_n,
\]
and if \(\ell(\varepsilon) \le \alpha_n\) the node is unsplit (only left child exists, same length).  All intervals are considered measurable subsets of \([0,\pi]\) with respect to Lebesgue measure.  The tree \(\mathcal{T}\) consists of all cylinder intervals.

\section{Generators from Split Nodes}
Let \(\varepsilon\) be a split node of depth \(n\) (so \(\ell(\varepsilon) > \alpha_n\)).  Its children are \(L = I_{\varepsilon0}\) and \(R = I_{\varepsilon1}\).  Define
\begin{equation}\label{eq:generator}
f_\varepsilon = \frac{\mathbf{1}_L}{|L|} - \frac{\mathbf{1}_R}{|R|}
= \frac{\mathbf{1}_{I_{\varepsilon0}}}{\ell(\varepsilon)-\alpha_n} - \frac{\mathbf{1}_{I_{\varepsilon1}}}{\alpha_n}.
\end{equation}
The function \(f_\varepsilon\) is supported on \(I_\varepsilon\), has zero mean, and its \(L^2\) norm is
\[
\|f_\varepsilon\|^2 = \frac{1}{|L|} + \frac{1}{|R|}
= \frac{1}{\ell(\varepsilon)-\alpha_n} + \frac{1}{\alpha_n}.
\]

\section{Gram–Schmidt Orthonormalisation}
The constant function \(\phi_0 := 1/\sqrt{\pi}\) is already normalised.  We order the split nodes by increasing depth, and for nodes at the same depth arbitrarily (e.g., from left to right).  Let \(\varepsilon^{(1)},\varepsilon^{(2)},\dots\) be this enumeration.  We build an orthonormal sequence \(\psi_1,\psi_2,\dots\) via the classical Gram–Schmidt process applied to \(\{f_{\varepsilon^{(k)}}\}\) in that order, always orthogonalising against all previously obtained wavelets and the constant.

\begin{lemma}\label{lem:gram}
The functions \(\psi_k\) obtained by Gram–Schmidt from \(\{f_{\varepsilon^{(k)}}\}\) form an orthonormal set.  Moreover, each \(\psi_k\) is a linear combination of \(f_{\varepsilon^{(j)}}\) with \(j\le k\) and is supported on a union of intervals from the tree.
\end{lemma}
\begin{proof}
Standard Gram–Schmidt preserves the linear span and, because the \(f_{\varepsilon}\) have finite support, produces an orthonormal system.  By induction, each \(\psi_k\) is a linear combination of the original generators up to index \(k\).
\end{proof}

We denote by \(\psi_\varepsilon\) the orthonormal function obtained from the node \(\varepsilon\) after this procedure; for the root we set \(\psi_\varnothing = \phi_0\).  The collection
\[
\mathcal{B} = \{\psi_\varnothing\} \cup \{\psi_\varepsilon : \varepsilon \text{ split node in }\mathcal{T}\}
\]
is then orthonormal by construction.

\section{Completeness}
\begin{theorem}\label{thm:complete}
The orthonormal system \(\mathcal{B}\) is complete in \(L^2[0,\pi]\).
\end{theorem}
\begin{proof}
Let \(V\) be the closed linear span of \(\mathcal{B}\).  We first observe that the original generators \(f_\varepsilon\) belong to \(V\) because they are finite linear combinations of the \(\psi_\eta\) with \(\eta\) preceding \(\varepsilon\) and \(\psi_\varepsilon\) itself (by the invertibility of the Gram–Schmidt matrix).  In particular, the constant function \(\mathbf{1}_{[0,\pi]}\) is a multiple of \(\psi_\varnothing\), so it lies in \(V\).

Now each \(f_\varepsilon\) equals a linear combination of the two indicators \(\mathbf{1}_L\) and \(\mathbf{1}_R\).  Conversely, the two indicators can be recovered from \(\mathbf{1}_{I_\varepsilon}\) and \(f_\varepsilon\) by solving a linear system:
\[
\mathbf{1}_L = \frac{|L|}{|I_\varepsilon|}\bigl( \mathbf{1}_{I_\varepsilon} + |R| f_\varepsilon \bigr),\qquad
\mathbf{1}_R = \frac{|R|}{|I_\varepsilon|}\bigl( \mathbf{1}_{I_\varepsilon} - |L| f_\varepsilon \bigr).
\]
Thus, starting from the constant function (which gives \(\mathbf{1}_{[0,\pi]}\)) and adding all \(f_\varepsilon\), we can successively generate the indicator functions of every cylinder interval.  Consequently, the space \(V\) contains the characteristic functions of all cylinder sets \(I_\varepsilon\) for all depths.

The collection of all cylinder intervals generates the Borel \(\sigma\)-algebra of \([0,\pi]\) and separates points (the tree is complete).  The linear span of their indicators is dense in \(L^2[0,\pi]\) by standard measure‑theoretic arguments (e.g., by approximating simple functions).  Since \(V\) contains this span, we have \(V = L^2[0,\pi]\).  Hence \(\mathcal{B}\) is an orthonormal basis.
\end{proof}

\section{The Harmonic Sine Transform}
We can now define the main object of the whole programme.

\begin{definition}[Harmonic Sine Transform]
For \(f\in L^2[0,\pi]\), its \textbf{Harmonic Sine Transform} (HST) is the map
\[
\mathcal{H}f : \mathcal{B} \to \mathbb{C},\qquad
(\mathcal{H}f)(\psi) = \langle f,\psi\rangle_{L^2[0,\pi]}.
\]
Because \(\mathcal{B}\) is an orthonormal basis, \(\mathcal{H}\) is an isometric isomorphism from \(L^2[0,\pi]\) onto the sequence space \(\ell^2(\mathcal{B})\).
\end{definition}

The inverse transform is given by the usual orthogonal expansion
\[
f = \sum_{\psi\in\mathcal{B}} \langle f,\psi\rangle \,\psi,
\]
with the sum converging unconditionally in \(L^2\).

\section{Conclusion}
Programme~II has been successfully executed.  We have constructed the greedy harmonic tree, extracted from it an orthonormal Haar basis via a combination of elementary zero‑mean generators and Gram–Schmidt orthogonalisation, and proved its completeness.  This basis enables the Harmonic Sine Transform, an isometric isomorphism that will be the lens through which we later view the transfer operator and the zeta function.

\begin{thebibliography}{9}
\bibitem{Geere2026a}
V.~Geere, \emph{A Purely Geometric Reconstruction of the Sine Function},
March 2026. \url{https://victorgeere.co.za/maths/sine-harmonic.html}
\bibitem{Geere2026b}
V.~Geere, \emph{On the Correlation Structure of a Greedy Harmonic Decomposition of the Sine Function},
March 2026. \url{https://victorgeere.co.za/maths/greedy-harmonic-decomposition.html}
\bibitem{ProgrammeI}
V.~Geere et al., \emph{Harmonic Sine Transform Programme I: Greedy Harmonic Expansion of Angles}, 2026.
\end{thebibliography}

\end{document}